\documentclass{ximera}
% xmPreamble.tex en xmPrintstyle.sty worden automatisch geladen door ximera.cls (v2.7.8);
% NIET nog eens \input{./preamble.tex} of \addPrintStyle{.} doen (geeft 'already defined'-fouten).

\makeatletter % eigen img-map toevoegen aan het zoekpad (werkt standalone én vanuit de xourse)
\g@addto@macro\Ginput@path{{img/}{UW4202/img/}}
\makeatother

\begin{document}
\author{Alexander Holvoet}
\xmtitle{Zand op een ellips}{}

\begin{exercise}
    Giet zoveel mogelijk droog zand op een ellips die in karton is uitgeknipt en horizontaal
    op een bekertje rust. Wat zie je?
    \begin{oplossing}
        In bovenaanzicht zien we een kamlijn: een stuk van de grote as, symmetrisch rond het
        middelpunt. Deze kamlijn lijkt niet tot aan de toppen op de grote as te gaan. Een
        eerste hypothese: misschien vallen de eindpunten van (het bovenaanzicht van) de
        kamlijn samen met de brandpunten van de ellips? In vooraanzicht zien we dat deze kam
        niet horizontaal loopt, maar in het midden hoger is dan opzij. Dat is niet
        verwonderlijk: de punten in het midden van de ellips liggen verder van de rand.
        \begin{image}[12cm]
            \includegraphics[width=6cm,keepaspectratio]{zand35.jpg}
            \quad
            \includegraphics[width=6cm,keepaspectratio]{zand36.jpg}
        \end{image}
    \end{oplossing}
\end{exercise}

Vanaf nu werken we in het bovenaanzicht. De punten van de kamlijn beschouwen we als punten
van de grote as van de ellips. We kunnen verwijzen naar de afrollende zandkorrels in de
ruimte, maar als we punten of rechten benoemen, doen we aan vlakke meetkunde.

\begin{exercise}
    Net als in de vorige opgaven kun je de kamlijn bekijken als een meetkundige plaats van
    punten die aan een bepaalde voorwaarde voldoen. Aan welke voorwaarde moeten de punten
    van de kamlijn voldoen?
    \begin{oplossing}
        De kamlijn is de meetkundige plaats van de punten op gelijke loodrechte afstand tot
        twee of meer randen van de figuur, hier dus tot de ellips. Zandkorrels rollen immers
        naar de dichtste rand; korrels aan de ene kant van de kamlijn rollen naar de ene
        rand, korrels aan de andere kant naar de andere rand.

        Voor de punten van de grote as \([DE]\) die niet te ver van het middelpunt liggen,
        klopt dit duidelijk. Punten van de grote as die dicht bij een top \(D\) of \(E\)
        liggen, rollen recht naar die top: zij liggen niet op een kamlijn, want er is maar
        \'e\'en dichtste randpunt. De vraag is dus: waar liggen de eindpunten \(A\) en \(B\)
        van de kamlijn, zo dat een punt \(P\) tussen \(A\) en \(B\) twee dichtste randpunten
        \(M\) en \(M'\) heeft, en een punt tussen \(A\) en \(D\) (of tussen \(B\) en \(E\))
        slechts \'e\'en, namelijk \(D\) (of \(E\))?
        \begin{image}[6cm]
            \includegraphics[width=6cm,keepaspectratio]{zand37.jpg}
        \end{image}
    \end{oplossing}
\end{exercise}

\begin{exercise}
    Je kunt het misschien omkeren. In plaats van een variabel punt \(P\) op de grote as te
    nemen en je af te vragen hoe dicht het bij de top moet liggen, kun je een variabel punt
    \(M\) van de ellips nemen, dat steeds dichter naar de top \(D\) laten gaan, en bepalen
    waar het bijbehorende punt \(P\) van de grote as in de limiet naartoe gaat. Onderzoek
    dit met GeoGebra voor de ellips \(\frac{x^2}{25} + \frac{y^2}{9} = 1\).
    \begin{oplossing}
        Teken deze ellips met GeoGebra, neem een punt \(M\) erop, en gebruik de knoppen
        `raaklijn' en `loodlijn' om de normaal in \(M\) te tekenen en met de grote as te
        snijden. Voor \(M\) heel dicht bij de top \(D\) heeft het punt \(P\) coördinaat
        \((3{,}2;\, 0)\). Dit is nog geen exacte methode, maar het lijkt er al op dat het
        eindpunt \(A\) van de kam niet het brandpunt is. Deze ellips heeft brandpunten
        \(F_{1,2}(\pm 4,\, 0)\).
        \begin{image}[12cm]
            \includegraphics[width=6cm,keepaspectratio]{zand38.jpg}
            \quad
            \includegraphics[width=6cm,keepaspectratio]{zand39.jpg}
        \end{image}
    \end{oplossing}
\end{exercise}

\begin{exercise}
    We willen nu het eindpunt \(A\) aan de rechterkant van de kam exact lokaliseren. Ken je
    een eigenschap van de ellips waar de normaal een rol in speelt?
    \begin{oplossing}
        De weerkaatsingseigenschap zegt dat \(F_1 M\) en \(F_2 M\) gelijke hoeken vormen met
        de normaal \(MP\) (waardoor een lichtstraal die in \(F_1\) vertrekt, in \(F_2\)
        toekomt en omgekeerd). Dus: \(MP\) is de bissectrice van de hoek
        \(\widehat{F_1 M F_2}\).
        \begin{image}[8.5cm]
            \includegraphics[width=8.5cm,keepaspectratio]{zand40.jpg}
        \end{image}
    \end{oplossing}
\end{exercise}

\begin{exercise}
    Wat zegt dit over de positie van punt \(P\) op de grote as? Preciezer: wat weet je over
    de verhouding \(\frac{|F_1 P|}{|F_2 P|}\)?
    \begin{hint}
        Pas de bissectricestelling toe in de driehoek \(F_1 M F_2\), met \(MP\) als
        bissectrice vanuit \(M\).
    \end{hint}
    \begin{oplossing}
        De bissectricestelling zegt: de bissectrice van een hoek van een driehoek verdeelt
        de overstaande zijde in stukken die in dezelfde verhouding staan als de aanliggende
        zijden. In de driehoek \(F_1 M F_2\) geeft dit
        \[
            \frac{|F_1 P|}{|F_2 P|} = \frac{|F_1 M|}{|F_2 M|} .
        \]
    \end{oplossing}
\end{exercise}

\begin{exercise}
    Maak de redenering af: waar komt \(P\) terecht als \(M\) naar de top \(D\) gaat? Maak
    het ten slotte concreet voor de ellips \(\frac{x^2}{25} + \frac{y^2}{9} = 1\).
    \begin{oplossing}
        Als \(M\) naar \(D\) gaat, komt \(P\) in het grenspunt \(A\) van de kamlijn
        \([AB]\). Het punt \(A\) is dus bepaald door
        \[
            \frac{|F_1 A|}{|F_2 A|} = \frac{|F_1 D|}{|F_2 D|} .
        \]
        Voor de ellips \(\frac{x^2}{25} + \frac{y^2}{9} = 1\) is de halve grote as \(5\) en
        de halve kleine as \(3\). We hebben \(D(5, 0)\), \(F_1(4, 0)\) en \(F_2(-4, 0)\).
        Stel \(A(u, 0)\); we zoeken \(u\) met
        \[
            \frac{4 - u}{u - (-4)} = \frac{5 - 4}{5 - (-4)} .
        \]
        Oplossen naar \(u\) geeft \(u = 3{,}2\), precies wat we met GeoGebra vonden. (We
        gingen ervan uit dat \(u < 4\), wat het GeoGebra-onderzoek overtuigend liet zien.)
    \end{oplossing}
\end{exercise}

\end{document}
